\documentclass{article}
\usepackage{amsmath}
\usepackage{amssymb}
\usepackage{amsthm}
\usepackage{MnSymbol}
\usepackage{xcolor}
\usepackage{parskip}
\usepackage{tikz}
\usepackage{bbm}
\usepackage{hyperref}
\usetikzlibrary{trees}
\usetikzlibrary{graphs}


\theoremstyle{definition}
\newtheorem{definition}{Definition}
\newtheorem{theorem}{Theorem}


\newcommand{\abs}[1]{\left| #1 \right|}


\title{lecture}
\date{\today}
\begin{document}
\maketitle

\section{betweeness (BTW)}
given a bunch of constraints ABC, b is between a and c, but dont care about order

random gives 1/3 since \(2/3!\) combinaions satisfy

assuming that the constrains are fully satisfied, we get a 1/2 approximation.


2 ways, sdp and combinatorial. 

\section{comb}

some of the vertecies are not between. (the farthest on either side cannot be in between since they are satisfiable.) these nodes are called ``free''.

place on either end. 

since there are two ends, the probability it matches any other constraints is 1/2

algo is linear

non-btw has random of 2/3 but has no better than 2/3 polytime algo

\section{SDP}


use whats called a ``feasability check'', then use embedding


theres no objective, what are the constraints?

inner product and distances are the same (one gives the answer to the other) with unit vectors

for constraint ijk, if we place them on a line, \(d_{ij}^2+d_{jk}^2\leq d_{ik}^2\). notice if the leq is =, then we get a triangle. given a random plane, the probability a projection of points that satisfy \(d_{n\choose 2}\) and also satisfy the original constraints is 1/2. since the inequality requires the triangle to be obtuse. 










\end{document}